\documentclass[../../../include/open-logic-section]{subfiles}

\begin{document}

\olfileid{sth}{spine}{zf}
\olsection{$\Z$ અને $\ZF$: એક મહત્ત્વનો પડાવ}

આધારથી આપણે બીજો એક મહત્ત્વનો પડાવ પાર કરીએ છીએ. અગાઉ
$\Zminus$ અને $\ZFminus$ વિશે કહ્યું હતું કે તે
અમુક બાબત ``બાદ કરેલા'' સિદ્ધાંતો છે. એ બાબત આધાર છે. તેથી:

\begin{defn}
સિદ્ધાંત $\Z$ એ $\Zminus$માં આધાર ઉમેરે છે. તેથી તેની
સ્વયંસિદ્ધિઓ આ છે: ઘટકો દ્વારા નિર્ધારિત સમાનતા, યોગગણ, જોડ,
ઘાતગણો, અનંતતા, આધાર અને પૃથક્કરણના પ્રરૂપના બધા કિસ્સાઓ.

સિદ્ધાંત $\ZF$ એ $\ZFminus$માં આધાર ઉમેરે છે. બીજા શબ્દોમાં,
$\ZF$ એ~$\Z$માં પ્રતિસ્થાપનના બધા કિસ્સાઓ ઉમેરે છે.
\end{defn}

તોપણ એક પ્રશ્ન થઈ શકે. $\ZFminus$ના સંદર્ભમાં નિયમિતતા
આધારને સમકક્ષ છે અને નિયમિતતાનું ઔચિત્ય સ્પષ્ટ છે, તો સીધી
નિયમિતતાને બદલે આધારને જ મૂળ સ્વયંસિદ્ધિ શા માટે લઈએ?

ઐતિહાસિક કારણો (કોણે શું અને ક્યારે ઘડ્યું) બાજુએ રાખીએ તો
મુખ્ય કારણ એ છે કે આધારને~$V_\alpha$ની વ્યાખ્યા વાપર્યા વિના
રજૂ કરી શકાય છે. એ વ્યાખ્યા
\olref[recursion]{sec}ના બધા કામ પર આધારિત હતી: તેને
વાજબી ઠેરવવા અનંતાતીત પુનરાવર્તન સાબિત કરવું પડ્યું.
પરંતુ આપણી અનંતાતીત પુનરાવર્તનની સાબિતીમાં
\emph{પ્રતિસ્થાપન} વપરાયું. તેથી આધાર અને નિયમિતતા
$\ZFminus$ના સંદર્ભમાં સમકક્ષ હોવા છતાં $\Zminus$ના
સંદર્ભમાં સમકક્ષ નથી.

વાસ્તવમાં બાબત આ ટૂંકી નોંધ કરતાં પણ વધુ ગંભીર છે.
આ પુસ્તકની મર્યાદાની બહારની વાત છે, પણ $\Zminus$ અને
$\Z$ બંને $V_\alpha$ વ્યાખ્યાયિત કરવા માટે પૂરતા શક્તિશાળી
નથી. તેથી માત્ર $\Z$માં કામ કરતાં હોય, તો આપણે જે રીતે
નિયમિતતા રજૂ કરી છે તે વિધાનનો \emph{અર્થ પણ નક્કી થઈ શકતો નથી}.
આ કારણે આપણી સત્તાવાર સ્વયંસિદ્ધિ નિયમિતતાને બદલે આધાર છે.

હવેથી જુદું ન કહીએ ત્યાં સુધી, વધુ ઉલ્લેખ કર્યા વિના,
આપણે $\ZF$માં કામ કરીશું.

\end{document}
